% !TeX root = main.tex
\documentclass[a4paper,12pt]{article}
% --------------------------------------------------
% Kodowanie i język
% --------------------------------------------------
\usepackage[T1]{fontenc}
\usepackage[utf8]{inputenc}
\usepackage[english,polish]{babel}
\usepackage{lmodern}

% --------------------------------------------------
% Strona
% --------------------------------------------------
\usepackage[a4paper,margin=2.5cm]{geometry}
\usepackage{setspace}
\onehalfspacing

% --------------------------------------------------
% Matematyka
% --------------------------------------------------
\usepackage{amsmath,amssymb,amsthm,mathtools}
\usepackage{bm}
\usepackage{nicefrac}

% --------------------------------------------------
% Grafika i tabele
% --------------------------------------------------
\usepackage{graphicx}
\usepackage{float}
\usepackage{caption}
\usepackage{subcaption}
\usepackage{booktabs}
\usepackage{array}
\usepackage{multirow}
\usepackage{tabularx}
\captionsetup{
    labelfont=bf,
    font=small
}

% --------------------------------------------------
% Kolory i linki
% --------------------------------------------------
\usepackage{xcolor}
\usepackage[hidelinks]{hyperref}
\usepackage[nameinlink,noabbrev]{cleveref}
\usepackage{csquotes}

% --------------------------------------------------
% Bibliografia
% --------------------------------------------------
\usepackage[
    backend=biber,
    style=numeric,
    sorting=none
]{biblatex}
\addbibresource{bibliography.bib}

% --------------------------------------------------
% LISTINGS (biały, clean styl)
% --------------------------------------------------
\usepackage{listings}

\definecolor{codebg}{RGB}{255,255,255}
\definecolor{codeframe}{RGB}{200,200,200}
\definecolor{codekw}{RGB}{0,0,180}
\definecolor{codecomment}{RGB}{0,120,0}
\definecolor{codestring}{RGB}{163,21,21}
\definecolor{codenum}{RGB}{120,120,120}
\definecolor{codetext}{RGB}{0,0,0}

\lstdefinelanguage{cppcustom}{
    language=C++,
    morekeywords={
        concept,requires,consteval,constinit,co_await,co_return,co_yield,
        constexpr,const_cast,dynamic_cast,static_cast,reinterpret_cast,
        nullptr,std,string,vector,map,unordered_map,set,unordered_set
    }
}

\lstdefinestyle{cppstyle}{
    language=cppcustom,
    backgroundcolor=\color{codebg},
    basicstyle=\ttfamily\small\color{codetext},
    keywordstyle=\bfseries\color{codekw},
    commentstyle=\itshape\color{codecomment},
    stringstyle=\color{codestring},
    numberstyle=\scriptsize\color{codenum},
    numbers=left,
    stepnumber=1,
    numbersep=10pt,
    showstringspaces=false,
    breaklines=true,
    breakatwhitespace=true,
    tabsize=4,
    frame=single,
    rulecolor=\color{codeframe},
    framesep=6pt,
    xleftmargin=12pt,
    framexleftmargin=8pt,
    keepspaces=true,
    columns=fullflexible,
    upquote=true
}

\lstset{style=cppstyle}

% --------------------------------------------------
% Twierdzenia
% --------------------------------------------------
\newtheorem{theorem}{Theorem}[section]
\newtheorem{lemma}[theorem]{Lemma}
\newtheorem{definition}[theorem]{Definition}
\newtheorem{example}[theorem]{Example}
\newtheorem*{remark}{Remark}


% --------------------------------------------------
% Makra matematyczne
% --------------------------------------------------
\newcommand{\R}{\mathbb{R}}
\newcommand{\N}{\mathbb{N}}
\newcommand{\Z}{\mathbb{Z}}
\newcommand{\E}{\mathbb{E}}
\newcommand{\Prob}{\mathbb{P}}

% --------------------------------------------------
% Makra tekstowe
% --------------------------------------------------
\newcommand{\cpp}{\texttt{C++}}
\newcommand{\code}[1]{\texttt{#1}}
\newcommand{\inlinecpp}[1]{\lstinline|#1|}
\newcommand{\pathname}[1]{\texttt{#1}}
\newcommand{\term}[1]{\emph{#1}}

% --------------------------------------------------
% Kod z pliku
% --------------------------------------------------
\newcommand{\cppfile}[2][]{
    \lstinputlisting[caption={Plik \pathname{#2}},#1]{#2}
}

% --------------------------------------------------
% Obrazki
% --------------------------------------------------
\newcommand{\img}[4][0.8\textwidth]{
    \begin{figure}[H]
        \centering
        \includegraphics[width=#1]{#2}
        \caption{#3}
        \label{#4}
    \end{figure}
}

\newcommand{\imgtwo}[5]{
    \begin{figure}[H]
        \centering
        \begin{subfigure}{0.48\textwidth}
            \centering
            \includegraphics[width=\textwidth]{#1}
            \caption{#2}
        \end{subfigure}
        \hfill
        \begin{subfigure}{0.48\textwidth}
            \centering
            \includegraphics[width=\textwidth]{#3}
            \caption{#4}
        \end{subfigure}
        \caption{#5}
    \end{figure}
}

% --------------------------------------------------
% Referencje
% --------------------------------------------------
\newcommand{\figref}[1]{\cref{#1}}
\newcommand{\tabref}[1]{\cref{#1}}
\newcommand{\secref}[1]{\cref{#1}}
\newcommand{\eqnref}[1]{\cref{#1}}

% --------------------------------------------------
% English names for cleveref
% --------------------------------------------------
\crefname{figure}{Fig.}{Figs.}
\Crefname{figure}{Fig.}{Figs.}
\crefname{table}{Table}{Tables}
\Crefname{table}{Table}{Tables}
\crefname{section}{Section}{Sections}
\Crefname{section}{Section}{Sections}
\crefname{equation}{Eq.}{Eqs.}
\Crefname{equation}{Eq.}{Eqs.}
